% ============================================================
%  ch1_intro  引言（原书 pp. 1--2，即 PDF 第 2--3 页）
% ============================================================

\section*{引言}
\addcontentsline{toc}{section}{引言}

\origpage{1}本文是 2001 年 2 月在 Tata 研究所讲授的十二小时课程（用英语讲授）的笔记；当时该所正举办以几何 Langlands 纲领为主题的特殊年。

作者无意把这些笔记整理成文，因此行文并不正式，毋宁说是想“讲述”函数域上 Langlands 对应之证明的总体脉络，并着重交代其中几个要点。作者希望这样做能便于读者进入那两篇体例完备却艰深的文章 \cite{lafforgue98,lafforgue02}。

第一章叙述 Langlands 对应本身，给出 Drinfeld chtoucas\footnote{译注：\emph{chtouca} 是 Drinfeld 造的词，本文一律保留原文、不作翻译，以免与“shtuka”“矢丛”等既有译名混淆。文献中常见的拼法还有 \emph{shtuka}，二者所指相同。[译者按（上下文推断，无外部依据）]} 及其水平结构的定义，讨论相应模叠的性质；接着介绍用 Harder--Narasimhan 典范多边形作截断的手法、不动点计数公式（取自 \cite{lafforgue97}），并预告紧化的构造，以及区分“本质上同调”与“可忽略上同调”的基本原理。

第二章给出无水平结构紧化的结构；陈述借助 Hecke 对应建立的“全局稳定性定理”，正是它使水平结构得以加回来；说明如何用配对 $L$-函数（自守侧与 Galois 侧都用）的论证把本质上同调与可忽略上同调分开；定义 Hecke 代数在本质上同调上的作用，并用迹来计算这个表示。所用的迹由一个“不动点定理”给出：该定理涉及的开集并不被 Hecke 对应整体稳定，只在“不动点附近局部地”稳定。

\origpage{2}第三章用“完整同态”构造无水平结构的紧化。更确切地说，作者引入一个“fourre-tout”\footnote{译注：\emph{fourre-tout} 法语原意为“什么都往里装的口袋”，此处指一个把所有截断紧化都收作开子叠的不分离叠。本文保留原文不译，理由见《译者说明》。[译者按（上下文推断，无外部依据）]}，它是一个不分离的叠，把对应于各个截断多边形的全部紧化都作为开子叠容纳其中；该章还给出边界分层的逐层分解（dévissage）。

第四章用截断在 fourre-tout 中定义各个开紧化，分析边界分层上诱导的截断，并陈述光滑性定理与紧合性主定理。该章陈述“限制到某个水平的光滑性定理”，并由此推出带水平结构的完全紧化或部分紧化的构造与主要性质。

第五章着手验证紧合性的赋值判据。它先聚焦于待退化的 chtouca 的一般纤维及其 $\varphi$-空间结构，引入该 $\varphi$-空间中“迭代 $\varphi$-格”的概念。该章说明：迭代 $\varphi$-格总有无穷多个，每一个都诱导 chtouca 的一个退化，而任务是要证明其中恰有一个具备所选截断多边形所规定的性质。为此引入“初等变换”的手法，用来从一个迭代 $\varphi$-格过渡到另一个。该章定义一般的“退化子”（dégénérateur）概念，追踪迭代 $\varphi$-格的初等变换中退化子族的走向，并由此推出“借助 Hecke 对应的全局稳定性定理”。

第六章（因时间不够未口头讲授）用一列初等变换构造出一个满足全部所需性质的迭代 $\varphi$-格。这一章比前面几章精细得多，并以略为不同的方式重述了 \cite{lafforgue98} 的证明；它还改正了该证明最后一页的一处（小）错误，这一错误由 Yakov Varshavsky 向作者指出（见文末的注记）。

作者十分感谢听课者的专注听讲与提问，感谢 Tata 研究所的数学家与负责人的邀请，也感谢 Yakov Varshavsky 的审读与细心，并再次感谢 IHES 的 Cheikhchoukh 女士：她一如既往，以极快的速度与细致的工夫完成了本文的打字工作。
